{\small\textbf{Banks (Zheng, 1:05) [7:50–8:55]}}\par\smallskip
Who would carry the losses if things went wrong? The banks, and they were in good shape.
\begin{itemize}\setlength\itemsep{0pt}
\item \textcolor{navy}{\textbf{[def]}} The core capital ratio, CET1, is high-quality equity as a share of risk-weighted assets. It was 17.7 percent at the end of 2021, above the EU average, and 16.3 percent a year later.
\item In DNB's stress test, the four major banks would lose 3.8 points of capital but stay well above the minimum.
\item Crucially, rising rates were \textbf{lifting} bank profits: net interest income grew 6.7 percent in 2022.
\end{itemize}
\par\smallskip
DNB noted that higher rates "may ultimately have a positive impact on the profitability of financial institutions". Keep that sentence in mind: it comes back as our twist.
\par\smallskip
So the setting cuts both ways. Against raising: war, inflation, rate hikes, housing turning. For raising: a recovered economy, cold credit, and strong, increasingly profitable banks.
\par\smallskip
\textcolor{nlred}{\textbf{HAND-OVER:}} "That was the setting. Diego will now show what DNB actually looked at, and what it was insuring against."
